\sect{पुरुषोत्तम-मासस्य शुक्ल-कमला-एकादशी-माहात्म्यम्}
\label{sec:padma-purushottama-shukla-kamala}

\ganapatyadiDhyanam
\padmaPuranam

\dnsub{अथ चतुःषष्टितमोऽध्यायः॥६४}

\uvacha{युधिष्ठिर उवाच}

\twolineshloka
{भगवञ्छ्रोतुमिच्छामि व्रतानामुत्तमं व्रतम्}
{सर्वपापहरं विष्णोः फलदं व्रतिनां च यत्}% ॥१॥

\twolineshloka
{पुरुषोत्तममासस्य कथां ब्रूहि जनार्दन}
{को विधिः किं फलं तस्य को देवस्तत्र पूज्यते}% ॥२॥

\twolineshloka
{अधिमासे च सम्प्राप्ते व्रतं ब्रूहि जनार्दन}
{कस्य दानस्य किं पुण्यं किं कर्तव्यं नृभिः प्रभो}% ॥३॥

\twolineshloka
{कथं स्नानं च किं जाप्यं कथं पूजाविधिः स्मृतः}
{किं भोज्यमुत्तमं चान्नं मासेऽस्मिन् पुरुषोत्तमे}% ॥४॥

\uvacha{श्रीकृष्ण उवाच}

\twolineshloka
{कथयिष्यामि राजेन्द्र भवतः स्नेहकारणात्}
{पुरुषोत्तममासस्य माहात्म्यं पापनाशनम्}% ॥५॥

\twolineshloka
{अधिमासे तु सम्प्राप्ते भवेदेकादशी तु या}
{कमला नाम सा नाम्ना तिथीनामुत्तमा तिथिः}% ॥६॥

\twolineshloka
{तस्या व्रतप्रभावेन कमलाभिमुखी भवेत्}
{ब्राह्मे मुहूर्ते चोत्थाय संस्मृत्य पुरुषोत्तमम्}% ॥७॥

\twolineshloka
{स्नात्वा चैव विधानेन नियमं कारयेद् व्रती}
{गृहे त्वेकगुणं जाप्यं नद्यां तु द्विगुणं स्मृतम्}% ॥८॥

\twolineshloka
{गवां गोष्ठे सहस्रोर्ध्वमग्न्यागारे शतान्वितम्}
{शिवक्षेत्रेषु तीर्थेषु देवतानां च सन्निधौ}% ॥९॥

\twolineshloka
{लक्षं तुलस्याः सान्निध्ये ह्यनन्तं विष्णुसन्निधौ}
{अवन्त्यामभवत् कश्चिच्छिवशर्मा द्विजोत्तमः}% ॥१०॥

\twolineshloka
{तस्याऽऽत्मजास्तु पञ्चाऽऽसन् कनिष्ठो दोषवानभूत्}
{तदा पित्रा परित्यक्तस्त्यक्तः स्वजनबान्धवैः}% ॥११॥

\twolineshloka
{कुकर्मणः प्रभावेन गतो दूरतरं वनम्}
{एकदा दैवयोगेन तीर्थराजं समागतः}% ॥१२॥

\twolineshloka
{क्षुत्क्षामो दीनवदनस्त्रिवेण्यां स्नानमाचरत्}
{मुनीनामाश्रमांस्तत्र विचिन्वन् क्षुधयाऽऽर्दितः}% ॥१३॥

\twolineshloka
{हरिमित्रमुनेस्तत्र ददर्शाश्रममुत्तमम्}
{पुरुषोत्तममासे वै जनानां च समागमे}% ॥१४॥

\twolineshloka
{तत्राऽऽश्रमे कथयतां कथां कल्मषनाशिनीम्}
{ब्राह्मणानां मुखात् तेन श्रद्धया कमला श्रुता}% ॥१५॥

\twolineshloka
{एकादशी पुण्यतमा भुक्तिमुक्तिप्रदायिनी}
{जयशर्मा विधानेन श्रुत्वेमां कमलातिथिम्}% ॥१६॥

\twolineshloka
{एकादशीपुण्यतमा भुक्तिमुक्तिप्रदायिनी}
{व्रतं कृतं च तैः सार्धं स्थित्वा शून्यालये तदा}% ॥१७॥

\twolineshloka
{निशीथे समनुप्राप्ते लक्ष्मीस्तत्र समागता}
{वरं ददामि भो विप्र कमलायाः प्रभावतः}% ॥१८॥

\uvacha{जयशर्मोवाच}

\twolineshloka
{का त्वं कस्यासि रम्भोरु प्रसन्ना च कथं मम}
{इन्द्राणी सुरनाथस्य भवानी शङ्करस्य वा}% ॥१९॥

\twolineshloka
{गन्धर्वी किन्नरी वाऽथ वधूर्वा चन्द्रसूर्ययोः}
{त्वत्सदृशा न दृष्टा च न श्रुता च शुभानने}% ॥२०॥

\uvacha{लक्ष्मीरुवाच}

\twolineshloka
{प्रसन्ना साम्प्रतं जाता वैकुण्ठादहमागता}
{प्रेरिता देवदेवेन कमलायाः प्रभावतः}% ॥२१॥

\twolineshloka
{पुरुषोत्तममासस्य या पक्षे प्रथमे भवेत्}
{तस्या व्रतं त्वया चीर्णं प्रयागे मुनिसन्निधौ}% ॥२२॥

\twolineshloka
{व्रतस्यास्य प्रभावेन वशगाहं न संशयः}
{तववंशे भविष्यन्ति मानवा द्विजसत्तम}% ॥२३॥

\onelineshloka*
{मत्प्रसादादवाप्स्यन्ति सत्यं ते व्याहृतं मया}

\uvacha{ब्राह्मण उवाच}
\onelineshloka
{प्रसन्ना यदि मे पद्मे व्रतं विस्तरतो वद}% ॥२४॥
\onelineshloka*
{यत्कथासु प्रवर्तन्ते साधवो ये जना द्विजाः}

\uvacha{लक्ष्मीरुवाच}
\onelineshloka
{श्रोतॄणां परमं श्राव्यं पवित्राणामनुत्तमम्}% ॥२५॥

\twolineshloka
{दुःस्वप्ननाशनं पुण्यं श्रोतव्यं यत्नतस्ततः}
{उत्तमः श्रद्धया युक्तः श्लोकं श्लोकार्धमेव वा}% ॥२६॥

\twolineshloka
{पठित्वा मुच्यते सद्यो महापातककोटिभिः}
{मासानां परमो मासः पक्षिणां गरुडो यथा}% ॥२७॥

\twolineshloka
{नदीनां च यथा गङ्गा तिथीनां द्वादशी तिथिः}
{अद्यापि निर्जराः सर्वे भारते जन्मलिप्सवः}% ॥२८॥

\twolineshloka
{तमर्चयन्ति विविधा नारायणमनामयम्}
{ये यजन्ति सदा भक्त्या देवं नारायणं प्रभुम्}% ॥२९॥

\twolineshloka
{तानर्चयन्ति सततं ब्रह्माद्या देवतागणाः}
{येऽपि नामपरा ये च हरिकीर्तनतत्पराः}% ॥३०॥

\twolineshloka
{हरिपूजापरा ये च ते कृतार्थाः कलौ युगेः}
{शुक्ले वा यदि वा कृष्णे भवेदेकादशीद्वयम्}% ॥३१॥

\threelineshloka
{गृहस्थानां भवेत् पूर्वा यतीनामुत्तरा स्मृता}
{एकादशी द्वादशी च रात्रिशेषे त्रयोदशी}
{तत्र क्रतुशतं पुण्यं त्रयोदश्यां तु पारणे}% ॥३२॥

\twolineshloka
{\textbf{एकादश्यां निराहारः स्थित्वाऽहमपरेऽहन्नि}}
{\textbf{भोक्ष्यामि पुण्डरीकाक्ष शरणं मे भवाच्युत}}% ॥३३॥

\twolineshloka
{अमुं मन्त्रं समुच्चार्य देवदेवस्य चक्रिणः}
{भक्तिभावेन तुष्टात्मा उपवासं समाचरेत्}% ॥३४॥

\twolineshloka
{कुर्याद् देवस्य पुरतो जागरं नियतो व्रती}
{गीतैर्वाद्यैश्च नृत्यैश्च पुराणपठनादिभिः}% ॥३५॥

\twolineshloka
{ततः प्रातः समुत्थाय द्वादशीदिवसे व्रती}
{स्नात्वा विष्णुं समभ्यर्च्य विधिवत् प्रयतेन्द्रियः}% ॥३६॥

\twolineshloka
{पञ्चामृतेन संस्नाप्य एकादश्यां जनार्दनः}
{द्वादश्यां च पयःस्नानाद्धरेः सारूप्यमश्नुते}% ॥३७॥

\twolineshloka
{अज्ञानतिमिरान्धस्य व्रतेनानेन केशव}
{प्रसीद सन्मुखो भूत्वा ज्ञानदृष्टिप्रदो भव}% ॥३८॥

\twolineshloka
{एवं विज्ञाप्य देवेशं देवदेवं गदाधरम्}
{ब्राह्मणान् भोजयेद् भक्त्या तेभ्यो दद्याच्च दक्षिणाम्}% ॥३९॥

\twolineshloka
{ततः स्वबन्धुभिः सार्धं नारायणपरायणः}
{कृत्वा पञ्चमहायज्ञान्स्वयं भुञ्जीत वाग्यतः}% ॥४०॥

\twolineshloka
{एवं यः प्रयतः कुर्यात् पुण्यमेकादशीव्रतम्}
{स याति विष्णुभवनं पुनरावृत्तिदुर्लभम्}% ॥४१॥

\twolineshloka
{इत्युक्त्वा कमला तस्मै वरं दत्त्वा तिरोदधे}
{सोऽपि विप्रो धनी भूत्वा पितुर्गेहं समागतः}% ॥४२॥

\uvacha{श्रीकृष्ण उवाच}

\twolineshloka
{एवं यः कुरुते राजन् कमलाव्रतमुत्तमम्}
{शृणुयाद् वासरे विष्णोः सर्वपापैः प्रमुच्यते}% ॥४३॥

॥इति श्रीपाद्मे महापुराणे पञ्चपञ्चाशत्साहस्र्यां संहितायामुत्तरखण्डे उमापतिनारदसंवादान्तर्गतकृष्णयुधिष्ठिरसंवादे पुरुषोत्तम-मासस्य शुक्ल-कमला-एकादशी-माहात्म्यं नाम चतुःषष्टितमोऽध्यायः॥६४॥

\genMangalaShloka

\hyperref[sec:ekadashi_mahatmyam_padma_puranam]{\closesub}
\clearpage

